\subsection{Grid Input}
\label{SEC:Files:Grid_Input}

%%%%%%%%%%%%%%%%%%%%%%%%%%%%%%%%%%%%%%%%%%%%%%%%%%%%%%%%%%%%%%%%%%%%%%%%%%%%%%%%

For specific test cases with trivial geometry, the flow solver will internally
generate a grid.
%
In general, though, the user must provide the mesh to be used for the
simulation.  
%
Grids that are two- or three-dimensional, structured or unstructured, can be
read from 
  CGNS,
  Plot3D, or
  Gmsh (very limited support)
files.
%
The GFR code will always assume the grid coordinates are in units of meters.
%
However, sometimes it can be easier or more useful to generate a grid in other
units~---~we refer to these as
\key[key:GridUnits]{\emph{grid units}}[grid units] which could be in anything
from millimeters, to inches, or even nondimensional~---~and the user-specified
input \inptext{Grid\_Scaling\_Factor} provides a way to scale these grid units
to meters as assumed by the solver.
%
See \sref{SEC:Grid_Scaling} and
\inp[nmlGrid::GridScalingFactor]{Grid\_Scaling\_Factor} for more information.

%%%%%%%%%%%%%%%%%%%%%%%%%%%%%%%%%%%%%%%%%%%%%%%%%%%%%%%%%%%%%%%%%%%%%%%%%%%%%%%%

If any of the surface geometries contain curvature, it is \emph{highly}
recommended to use a high-order curved grid with at least \PO[2] grid cells to
accurately represent the surfaces with curvature.
%
This recommendation is simply so that the grid represents the curved boundary
surfaces with at least some amount of curvature and not linear segments.
%
Increasing the grid order to \PO[3+] can help improve the simulation accuracy,
but how much of an improvement depends on the flow conditions, the amount of
curvature in individual cells (think coarse cells with high curvature), etc.
%
One example scenario that can show improved accuracy and stability from
increasing the grid order to \PO[3+] is when significant regions of the geometry
surfaces are defined by circles, \thatis cylindrical\slash axisymmetric or
spherical geometries, simply because a \PO[3+] polynomial provides a markedly
better approximation of a circle compared to a \PO[2] polynomial approximation.
%
Note, the grid order and solution order are completely independent.
%
The grid order does not need to match or exceed the solution order, or vice
versa.

%%%%%%%%%%%%%%%%%%%%%%%%%%%%%%%%%%%%%%%%%%%%%%%%%%%%%%%%%%%%%%%%%%%%%%%%%%%%%%%%

Grid files are currently subject to the following limitations:
%
\begin{itemize}
  \item
    GFR does not support structured zones in CGNS grid files.
  \item
    Gmsh grid files should be avoided since GFR only supports Gmsh files in the
    older ASCII Legacy Version 2 format, and the Gmsh grid reader has been
    minimally tested.
  \item
    Multizone grids must have contiguous, point-matched, non-overlapping faces.
  \item
   %Only 2D quadrilateral or 3d hexahedral elements are permitted.
    Only the element types described in \sref{SEC:Element_Types} are permitted.
  \item
    GFR does not support elements with collapsed faces or edges. 
    %
    For example, this can happen if trying to create a 3D axisymmetric
    unstructured grid by rotating a 2D unstructured grid about a pole boundary,
    and the 2D grid contains one or more quad cells with only a single node on
    the pole boundary.
    %
    This will create an invalid hexahedral element with a single collapsed edge
    at the pole boundary.
    %
    In this particular situation, it is generally recommended to create the 2D
    grid such that it does not go all the way to the pole boundary \slash
    rotation axis, create the 3D grid by rotating the 2D grid about the rotation
    axis, and finally create a separate 3D grid within the void that is created
    along the rotation axis.
\end{itemize}

%%%%%%%%%%%%%%%%%%%%%%%%%%%%%%%%%%%%%%%%%%%%%%%%%%%%%%%%%%%%%%%%%%%%%%%%%%%%%%%%

The flow solver hides much of the complicated grid manipulation from the user.
%
During the flow solution process, structured grids are converted to unstructured
quadrilaterals or hexahedra, multizone grids are merged, and redundant faces are
removed.
%
The required number of solution points will then be inserted into the cell
volume for the specified solution order.
%
The coordinates of these points are determined using linear interpolation
between the cell vertices and the chosen Gauss or Lobatto locations.
%
The domain is then partitioned to make efficient use of the computational
resources available.
%
Before output files are generated, the domain is recombined to simplify flow
visualization.

%%%%%%%%%%%%%%%%%%%%%%%%%%%%%%%%%%%%%%%%%%%%%%%%%%%%%%%%%%%%%%%%%%%%%%%%%%%%%%%%

\SeeAlso{
  \inp[nmlGrid::GridFile]{Grid\_File},
  \inp[nmlGrid::GridFormat]{Grid\_Format},
  \inp[nmlGrid::GridScalingFactor]{Grid\_Scaling\_Factor},
  \inp[nmlInitialization::Initialization]{Initialization},
  \inp[nmlNumerics::LocSolutionPts]{Loc\_Solution\_Pts},
  \inp[nmlNumerics::SolutionOrder]{Solution\_Order}
}

%%%%%%%%%%%%%%%%%%%%%%%%%%%%%%%%%%%%%%%%%%%%%%%%%%%%%%%%%%%%%%%%%%%%%%%%%%%%%%%%

